\begin{abstract}
Reference-guided prompting is attractive when dense medical annotations are scarce: a frozen visual matcher can localize evidence in a query image, and a frozen segmentation foundation model can in principle refine that evidence into a mask. This Phase 3 report tests that premise on the GlaS gland segmentation protocol after a pre-registered Phase 2D failure of grouped sparse point prompts. The original operational state, \texttt{STOP\_GLAS\_SEGMENTER\_FAILURE\_CONFIRMED}, is preserved, but the scientific interpretation is narrowed: the evidence supports failure of the grouped sparse point interface rather than complete failure of the SAM backbone. Privileged GT box prompts remain effective, motivating a bounded automatic region-prompt feasibility study. We pre-register a fixed region proposal bank, freeze DINOv2-L matching evidence, split the 55-image utility pool into 40/15 internal pools, and evaluate Gate R1 on 30 strict-validation queries before any new SAM region inference. No automatic region family passes Gate R1. The best automatic box route, R03 automatic union box, reaches mean box IoU 0.745 but improves over size-matched random boxes by only 0.026 with win rate 0.200, far below the pre-registered thresholds. The automatic route therefore stops before SAM refinement, oracle analysis, router training, or testA/testB access. This negative result is reported as a gated boundary condition for prompt representation selection, not as an automatic region prompting success.
\end{abstract}
